\documentclass[11pt]{article}
% Shared preamble for the three living documents.
\usepackage[T1]{fontenc}
\usepackage{lmodern}
\usepackage[utf8]{inputenc}
\usepackage[margin=2.4cm]{geometry}
\usepackage{microtype}
\usepackage{amsmath,amssymb}
\usepackage{siunitx}
\usepackage{booktabs}
\usepackage{enumitem}
\usepackage{xcolor}
\usepackage{tikz}
\usetikzlibrary{arrows.meta,positioning,calc}
\usepackage{listings}
\usepackage[backend=biber,style=numeric,sorting=none]{biblatex}
\addbibresource{references.bib}
\usepackage[colorlinks=true,linkcolor=blue!50!black,citecolor=green!40!black,urlcolor=blue!60!black]{hyperref}

\newcommand{\vect}[1]{\mathbf{#1}}
\newcommand{\unitv}[1]{\hat{\mathbf{#1}}}
\newcommand{\dd}{\mathrm{d}}
\newcommand{\project}{Physics Sandbox}

\lstset{basicstyle=\ttfamily\small,columns=flexible,keepspaces=true,breaklines=true,frame=single,framerule=0.3pt}
\setlist{itemsep=2pt}
\setlength{\emergencystretch}{3em}

\title{\project{} --- Physics Notes}
\author{Living document}
\date{Last updated: \today}

\begin{document}
\maketitle
\noindent This document records the physical models used by the sandbox, the
assumptions behind them and their known limitations. Each implementation step
adds a section. The companion documents cover the numerical methods
(\emph{Math Notes}) and the software design (\emph{Code Notes}).

\tableofcontents

\section{Conventions}
\begin{itemize}
  \item SI units throughout: metres, kilograms, seconds, newtons.
  \item Right-handed coordinates with $y$ pointing \emph{up} (the three.js
        convention~\cite{threejs}). The floor is the plane $y=0$; the table top is
        at $y=h=\SI{0.8}{m}$.
  \item A force direction is given by an \emph{azimuth} $\theta$ in the horizontal
        $x$--$z$ plane, measured from $+x$ toward $-z$ (counter-clockwise when seen
        from above), and an \emph{elevation} $\varphi$ above the horizontal:
        \begin{equation}
          \unitv{n}(\theta,\varphi) = (\cos\varphi\cos\theta,\; \sin\varphi,\; -\cos\varphi\sin\theta).
          \label{eq:dir}
        \end{equation}
\end{itemize}

\section{Milestone 1: a ball on a table}

\subsection{Model}
The ball (mass $m$, radius $r$) is treated as a \emph{particle} located at its
centre $\vect{r}(t)$. Its motion obeys Newton's second law~\cite{halliday2013,kleppner2014}
\begin{equation}
  m\,\ddot{\vect{r}} = \vect{F}_\text{net} = \vect{F}_g + \sum_i \vect{F}_i(t) + \vect{N} + \vect{f}.
\end{equation}
The radius matters only for \emph{where} contact happens: the ball touches a
surface of height $y_s$ when its centre is at $y = y_s + r$.

\subsection{Forces}
\paragraph{Gravity.} $\vect{F}_g = (0,-mg,0)$ with $g = \SI{9.81}{m/s^2}$ by default.

\paragraph{Applied forces.} The user specifies a magnitude $F_i$, a direction
\eqref{eq:dir}, and a duration $\tau_i$. The force starts at the time $t_i$ at
which it is applied and is constant on the window $t_i \le t < t_i+\tau_i$:
\begin{equation}
  \vect{F}_i(t) = F_i\,\unitv{n}_i \quad\text{if } t\in[t_i, t_i+\tau_i),\qquad \vect{0}\text{ otherwise.}
\end{equation}
The impulse delivered is $\vect{J}_i = F_i\tau_i\,\unitv{n}_i$, so, on a frictionless
table, a force acting for $\tau$ changes the velocity by $\Delta\vect{v} = \vect{J}/m$.

\paragraph{Normal force.} A rigid horizontal surface can only push. With
$\vect{F}_\text{ext} = \vect{F}_g + \sum_i\vect{F}_i$, a ball resting on the surface feels
\begin{equation}
  N = \max\!\left(0,\, -F_{\text{ext},y}\right), \qquad \vect{N} = (0, N, 0).
\end{equation}
If the applied force pulls upward harder than the weight ($F_{\text{ext},y}>0$) the
normal force vanishes and the ball lifts off.

\paragraph{Friction (Coulomb model).} The classic dry-friction law~\cite{halliday2013,stewart2000}:
\begin{itemize}
  \item \emph{Static} (ball at rest, horizontal push $\vect{F}_h$): $\vect{f} = -\vect{F}_h$ as long as
        $|\vect{F}_h| \le \mu_s N$; the ball does not move.
  \item \emph{Kinetic} (ball sliding with horizontal velocity $\vect{v}_h$):
        $\vect{f} = -\mu_k N\,\vect{v}_h/|\vect{v}_h|$, independent of speed.
  \item If a resting ball is pushed with $|\vect{F}_h| > \mu_s N$ it breaks free and
        kinetic friction acts opposite to the push.
\end{itemize}
Default values: table $\mu_s = 0.25$, $\mu_k = 0.20$; floor $\mu_s = 0.5$, $\mu_k=0.4$.

\subsection{Leaving the table and landing}
The table supports the ball only while its centre lies above the table top
($|x|\le \SI{1.2}{m}$, $|z|\le\SI{0.6}{m}$). Past the edge the ball is a projectile: horizontal
velocity is constant and it accelerates downward at $g$~\cite{taylor2005}. Starting from the
table with horizontal speed $v$, it reaches the floor after
\begin{equation}
  t_\text{fall} = \sqrt{2h/g} \approx \SI{0.404}{s}, \qquad \text{landing a horizontal distance } v\,t_\text{fall}\text{ from the edge.}
\end{equation}
Landing is modelled as \emph{perfectly inelastic} in the vertical direction (no bounce): the
vertical velocity is removed and the ball then slides on the floor with the floor's friction.

\subsection{Reference results (used as automated tests)}
\begin{center}
\begin{tabular}{@{}ll@{}}
\toprule
Situation & Expected result \\ \midrule
Constant $F$ on frictionless table & $x = \tfrac12 (F/m)\, t^2$, $v = (F/m)\,t$ \\
$F \le \mu_s m g$ & ball stays at rest \\
$F > \mu_s m g$ & $a = (F - \mu_k m g)/m$ \\
Sliding at $v_0$, no applied force & stops after $t = v_0/(\mu_k g)$, distance $v_0^2/(2\mu_k g)$ \\
Off the edge & reaches the floor after $\sqrt{2h/g}$ \\
Upward force $> mg$ & ball lifts off \\ \bottomrule
\end{tabular}
\end{center}

\subsection{Kinematic quantities shown in the plots}
\begin{itemize}
  \item \textbf{Displacement} $\Delta\vect{r}(t) = \vect{r}(t)-\vect{r}(0)$: a vector; components $\Delta x,\Delta y,\Delta z$ and magnitude $|\Delta\vect{r}|$.
  \item \textbf{Distance} (path length) $s(t) = \int_0^t |\vect{v}|\,\dd t'$: a scalar that never decreases.
        The ``components'' plotted are the per-axis path lengths $\int_0^t |v_x|\,\dd t'$ etc.
        Note $s \ge |\Delta\vect{r}|$, with equality only for straight-line motion without reversal.
  \item \textbf{Velocity} $\vect{v} = \dot{\vect{r}}$ and \textbf{speed} $|\vect{v}|$.
  \item \textbf{Acceleration} $\vect{a} = \vect{F}_\text{net}/m$ and its magnitude.
\end{itemize}

\subsection{Assumptions and limitations}
\begin{itemize}
  \item \textbf{No rotation.} A real ball pushed on a rough table would roll. For a uniform
        solid sphere rolling without slipping, $I = \tfrac25 m r^2$ and a horizontal force through
        the centre gives $a = \tfrac{5}{7}F/m$~\cite{kleppner2014}. Rolling, and the transition
        from sliding to rolling, is planned for a later milestone.
  \item No air resistance.
  \item The ball is supported while its \emph{centre} is over the table; it does not tip
        over the edge gradually, and collisions with the side of the table are ignored.
  \item Landing is perfectly inelastic (coefficient of restitution $e=0$).
  \item Table and floor are rigid and immovable.
\end{itemize}

\section{Milestone 1.1: configurable table}
\subsection{Table geometry}
The table top is a horizontal rectangle of length $L$ (along $x$), width $W$ (along $z$) and height $h$,
centred above the origin. The user chooses all three
($\SI{0.2}{m}\le L, W\le\SI{20}{m}$, $\SI{0.1}{m}\le h\le\SI{3}{m}$). The physics needs only the
supporting region $|x|\le L/2$, $|z|\le W/2$ at $y=h$. The ball starts at rest at
$(-L/3,\ h+r,\ 0)$, on the left third of the table, which leaves room to push it toward $+x$.

Changing $h$ changes how long the ball takes to fall after leaving the edge,
$t_\text{fall}=\sqrt{2h/g}$, and so how far from the table it lands. This makes for an easy
classroom experiment: double the height and the landing distance grows by $\sqrt2$.

\subsection{No table}
The table can be removed. The ball then starts at rest on the floor at $(0, r, 0)$, and every
force acts against the floor's friction coefficients. The floor is unbounded, so the ball can never
fall off.

\subsection{Choosing a force by dragging}
The force can also be set by dragging the tip of the preview arrow in the 3D view. The arrow is
drawn to a fixed scale of \SI{0.06}{m} per newton, so the arrow's length \emph{is} the
magnitude. Dragging changes the force vector directly, and the magnitude, azimuth and
elevation fields update to match (see the \emph{Math Notes}).

\section{Milestone 1.2: reading the forces}
Hovering over any force arrow (or tapping it on a touch screen) shows that force's magnitude,
its Cartesian components and its direction as azimuth and elevation. Because the arrows are
re-evaluated at the current time, this makes a free-body diagram that can be read off directly.
For a ball resting on the table, for instance, gravity reads $(0, -mg, 0)$, the normal force reads
$(0, +mg, 0)$, and their sum is zero, as Newton's first law requires~\cite{halliday2013}. While an
applied force pushes the ball, the friction arrow switches from \emph{static} (equal and opposite to
the push) to \emph{kinetic} (magnitude $\mu_k N$, opposite to the velocity).

The applied-force arrow shows the \emph{sum} of all applied forces active at that moment.

\section{Milestone 2: bouncing, energy, and an editable scene}

\subsection{Bouncing: the coefficient of restitution}
In a real impact the ball compresses, then springs back, and loses some energy as heat, sound and
vibration. This is summarised by the \emph{coefficient of restitution}~\cite{halliday2013,taylor2005}
\begin{equation}
  e = \frac{\text{speed of separation}}{\text{speed of approach}} = -\frac{v_y^{\,\text{after}}}{v_y^{\,\text{before}}},\qquad 0\le e\le 1.
\end{equation}
$e=0$ means the ball does not bounce (the earlier behaviour), and $e=1$ means a perfectly elastic bounce. Since kinetic energy goes as
$v^2$, each bounce keeps a fraction $e^2$ of the vertical kinetic energy, so a ball dropped from height $h$ rebounds to
\begin{equation}
  h_1 = e^2 h,\qquad h_n = e^{2n} h .
\end{equation}
Restitution really belongs to a \emph{pair} of materials. The sandbox gives the ball and each surface their own value and uses
$e = e_\text{ball}\, e_\text{surface}$ for an impact. Surfaces default to $1$, so the ball's value applies unless a surface is
made ``softer''. Only the vertical (normal) velocity is affected; friction during the brief impact is ignored. A bounce slower than
$2g\,\Delta t$ is treated as coming to rest, which avoids an endless series of ever-tinier bounces.

\subsection{Mechanical energy}
The energy plot shows
\begin{equation}
  \text{KE} = \tfrac12 m |\vect{v}|^2,\qquad \text{PE} = m g\,(y - r),\qquad \text{TE} = \text{KE} + \text{PE}.
\end{equation}
Potential energy is measured from the ball resting on the floor (centre at $y=r$). Only differences in PE matter
physically, so the choice of zero is a convention~\cite{kleppner2014}. The \emph{work--energy theorem} says
\begin{equation}
  \Delta\text{TE} = W_\text{applied} + W_\text{friction} + \Delta E_\text{impacts},
\end{equation}
so the plot shows directly:
\begin{itemize}
  \item free fall: KE and PE exchange while TE stays constant (gravity is conservative);
  \item an applied force: TE rises by the work $\int \vect{F}\cdot\vect{v}\,\dd t$;
  \item sliding friction: TE falls steadily (energy converted to heat);
  \item each bounce: a sudden drop of TE by the factor $(1-e^2)$ of the vertical kinetic energy.
\end{itemize}

\subsection{Initial conditions and editing}
Every object in the scene can be selected and edited:
\begin{itemize}
  \item the ball's mass, radius, bounciness, starting position and starting velocity;
  \item the table's size, height, position, friction and bounciness (or the table can be removed);
  \item the floor's friction and bounciness;
  \item gravity (for instance Earth \SI{9.81}{m/s^2}, Moon \SI{1.62}{m/s^2}, Mars \SI{3.71}{m/s^2}).
\end{itemize}
These are \emph{initial conditions} and \emph{parameters}. Newtonian motion is fully determined by them~\cite{taylor2005},
so every edit re-runs the simulation from $t=0$. Two rules keep the scene physically sensible: the ball cannot be placed inside
the table top or below the floor, and a ball resting on the table moves with the table when the table is moved or resized.

\section*{Changelog}
\begin{itemize}
  \item \textbf{Milestone 1} --- particle on a table: gravity, timed applied forces, normal force,
        Coulomb friction, edge fall, inelastic landing.
  \item \textbf{Milestone 1.1} --- user-defined table size and height, optional table, force
        chosen by dragging its arrow.
  \item \textbf{Milestone 1.2} --- force tooltips (free-body diagram readout).
  \item \textbf{Milestone 2} --- coefficient of restitution (bouncing), energy plot, editable initial
        conditions (position, velocity) and object properties.
\end{itemize}

\printbibliography
\end{document}
